\documentclass[11pt]{article}
\usepackage[T1]{fontenc}
\usepackage[utf8]{inputenc}
\usepackage[english]{babel}
\usepackage[margin=1in]{geometry}
\usepackage{amsmath,amssymb}
\usepackage{enumitem}
\usepackage{microtype}

\ifdefined\pdfcatalog
  \pdfcatalog{/Lang (en-US)}
\fi

\setlength{\parindent}{0pt}
\setlength{\parskip}{0.55em}
\setlength{\emergencystretch}{2em}
\setlist{nosep,leftmargin=*}

\newcommand{\findinglabel}[1]{\textbf{#1}}
\newcommand{\classification}[1]{\textit{Classification: #1.}}

\title{Technical Review of ``Early-Stopping Activation Steering for Factual Preference Alignment in Vietnamese Medical Language Models''}
\author{Manuscript review}
\date{August 24, 2026}

\begin{document}
\selectlanguage{english}
\maketitle

\section*{Overall assessment}

The manuscript has made several substantive corrections.  All headline conditions now use the same 500-question denominator, the bounded-generation McNemar calculation is labeled correctly, negative steering and 20 random directions are reported, the BM25 condition is identified as non-oracle, and the source compiles.  The strongest supported result is therefore clearer: under the 200-token bounded-generation protocol, early-stop steering increases the manuscript's BERTScore reference-preference proxy from 365/500 to 386/500, with a valid two-sided exact McNemar result for the displayed contingency table.

The paper is not yet ready for acceptance, however.  The primary natural-completion result is small and nonsignificant; the manuscript does not directly show that early stopping prevents long-sequence saturation; the schedule ablation lacks uncertainty and leaves dose-shape confounds; the placebo summary treats a standardized distance as stronger evidence than 20 random draws justify; and the BM25 comparison is not sufficiently specified or statistically compared.  Dataset construction, clinical validation, vector extraction, hook implementation, and systems timing also remain incomplete.  Major revision is still required, with emphasis on the mechanism and clinically meaningful evaluation rather than additional headline language.

\section*{Major technical findings}

\subsection*{1. The experiments do not isolate the claimed early-stopping mechanism}

\classification{Unsupported central mechanism claim}

\findinglabel{Location.} Abstract; Introduction; Early-Stopping Steering Hook; Controlled Factorial Ablation Matrix; Discussion; Conclusion.

\findinglabel{Problem.} The manuscript states that continuous steering causes ``representation saturation'' over long sequences and presents early stopping as the mechanism that prevents it.  No activation-norm, directional-drift, saturation, or hidden-state measurement is reported.  More importantly, the primary 800-token experiment compares only the unsteered baseline with early-stop steering; it does not include continuous steering, hard cutoff, or a dose-matched control.  The 200-token ablation provides only point estimates.  Its repetition results are not uniformly favorable to early stopping: at $\alpha_0=15$, continuous steering has lower Rep-4 than both hard cutoff and linear decay (3.65\% versus 3.73\% and 4.03\%).  At $\alpha_0=18$ and 20, the reported advantages are small and have no confidence intervals.

\findinglabel{Why it matters.} The results can support an empirical effect for one early-stop configuration, but they do not establish that saturation occurs or that stopping/decay is the cause of the effect.  The paper's methodological novelty depends on this distinction.

\findinglabel{Suggested fix.} Under naturally terminated generation, compare continuous steering, hard cutoff, linear decay, an unsteered baseline, equal-$D_1$ controls, and controls matched as closely as possible on $D_2=\sum_t\alpha(t)^2$.  Report paired differences and confidence intervals for reference preference, BERTScore, Rep-4, output length, and clinically adjudicated errors.  Add activation-level diagnostics if ``representation saturation'' is retained as a mechanism.  Otherwise describe saturation as a hypothesis and present early stopping as a temporal steering schedule rather than a demonstrated prevention mechanism.

\subsection*{2. The central outcome remains a semantic proxy rather than factual or clinical correctness}

\classification{Unsupported interpretation of a correctly qualified proxy}

\findinglabel{Location.} Abstract; Introduction; contributions; Results; Conclusion; keywords.

\findinglabel{Problem.} The Methods section appropriately states that RefPref is a BERTScore proxy and not clinician-adjudicated correctness.  Elsewhere, however, the manuscript repeatedly converts that proxy into ``factual alignment,'' ``truthfulness,'' ``superior accuracy,'' and hallucination mitigation.  The BM25 comparison is also described as an 8.60-point factual-alignment advantage even though it is an 8.60-point difference in reference-preference rate.  A response can be closer to $y^+$ than to one synthetic $y^-$ while retaining a wrong dose, unit, negation, contraindication, or interaction mechanism.

\findinglabel{Why it matters.} The distinction is especially important in clinical text, where small semantic or numeric errors can be harmful and may not be detected reliably by embedding similarity.

\findinglabel{Suggested fix.} Make blinded clinician-adjudicated factual correctness and harmful-error rate the primary endpoints, stratified by pregnancy, interaction, and dosage questions.  Validate RefPref against those labels and report sensitivity, specificity, and failure cases.  Until then, use ``reference-preference rate'' consistently in the Abstract, Results, and Conclusion; avoid ``clinical accuracy,'' ``truthfulness,'' and hallucination-mitigation claims as measured outcomes.

\subsection*{3. The primary natural-completion experiment is incomplete and does not support an improvement claim}

\classification{Definite reporting omission and unsupported behavioral inference}

\findinglabel{Location.} Primary Natural Completion Evaluation; Discussion; Abstract; Conclusion.

\findinglabel{Problem.} The reported change is 324/500 to 329/500 with $p=0.4583$.  A net difference of five does not reproduce a paired test; the discordant counts must be given.  The stated exact binomial value is compatible, for example, with 12 regressions and 17 improvements, but the manuscript does not report that table or its test definition.  The BERTScore change $0.6804\rightarrow0.6834$ has no paired interval or test.  Rep-4, output-length distribution, finish reasons, and category effects are omitted.  A 100\% EOS hit rate establishes that all outputs emitted EOS before the cap, but it does not establish that they contained no repetition loops before termination.

\findinglabel{Why it matters.} This is labeled the primary, deployable protocol, yet it does not show a statistically supported preference gain or document the fluency behavior that motivates early stopping.  The larger 4.20-point effect occurs only in the bounded stress test.

\findinglabel{Suggested fix.} Report the complete paired transition table, exact test, paired BERTScore effect and confidence interval, output-length distribution, Rep-4, finish reasons, and category-specific estimates.  Include continuous steering in this protocol.  Describe the present result as a nonsignificant trend and remove ``without repetition loops'' unless repetition is measured directly.

\subsection*{4. The multi-random placebo analysis is useful but statistically overstated}

\classification{Definite aggregation error and unsupported strength of evidence}

\findinglabel{Location.} Abstract; contribution~3; Contrastive Vector Estimation \& Directional Controls; Placebo Control Distribution; Table~\texttt{tab:baselines}; Conclusion.

\findinglabel{Problem.} The 20-direction experiment is a meaningful improvement over a single placebo vector, but the reporting remains incomplete.  The quantity
\[
\frac{77.20-56.50}{5.11}\approx4.05
\]
is a standardized distance from the sample mean, not a calibrated hypothesis test or confidence level.  With only 20 sampled random directions and none exceeding the target, the simplest finite-sample randomization statement is an empirical one-sided bound of approximately $1/(20+1)=0.0476$, subject to a prespecified sampling procedure.  The table's ``Raw Count'' of 283/500 is invalid for a distributional mean: 56.50\% corresponds to 282.5 successes per 500 on average, and the experiment actually contains 20 sets of 500 outputs.  BERTScore F1 is shown as a single value (0.6945) without stating whether it is a mean or reporting dispersion.  Seeds, individual direction results, and paired uncertainty are absent.  Negative steering is substantially asymmetric relative to positive steering ($-10.20$ versus $+4.20$ points), but the text describes it as a symmetric-direction test without analyzing that asymmetry.

\findinglabel{Why it matters.} The controls support direction-specific evidence, but they do not by themselves prove that the vector encodes factuality or establish a $4\sigma$ causal result.  Ambiguous aggregation prevents reproduction.

\findinglabel{Suggested fix.} Release all 20 seeds and per-direction, per-question outputs.  Replace the raw-count cell with mean $\pm$ SD and range across directions for every metric.  Use a prespecified randomization/permutation analysis and a hierarchical or nested bootstrap over both questions and sampled directions.  Compare $+v_{\text{steer}}$, $-v_{\text{steer}}$, and random directions on the identical paired records.  State that the result is evidence of directional specificity, not proof of a factual representation.

\subsection*{5. The BM25 RAG comparison is not reproducible or sufficient for a superiority claim}

\classification{Likely weak baseline and unsupported causal attribution}

\findinglabel{Location.} Experimental Setup; Placebo Control Distribution \& Real BM25 RAG Results; Abstract; Conclusion.

\findinglabel{Problem.} The paper reports Top-1 Recall@1 of only 13.00\% over 4,495 passages, but it does not define the corpus construction, passage boundaries, Vietnamese tokenization, BM25 implementation and parameters, query text, relevance mapping, prompt template, or handling of retrieval failures.  It attributes the 68.60\% reference-preference result ``to retrieval noise'' without an oracle-context control or error decomposition.  The steering-versus-RAG difference has no paired contingency table or confidence interval.  Calling this condition ``realistic'' and concluding that steering is superior is therefore premature; a 13\% retrieval recall may instead indicate an underdeveloped retriever.

\findinglabel{Why it matters.} The systems and accuracy comparison is a headline contribution.  A weak or irreproducible RAG implementation can exaggerate both steering's apparent quality advantage and the practical trade-off.

\findinglabel{Suggested fix.} Fully specify and release the retrieval pipeline.  Report Recall@$k$ for several $k$, a gold-context/oracle condition, retrieval-only error analysis, and generation results conditioned on successful versus failed retrieval.  Compare methods on the same questions using paired intervals/tests.  Separate retrieval time, prompt-prefill time, decoding time, index memory, and KV-cache memory.  Describe the current result as one BM25 configuration, not a general comparison with RAG.

\subsection*{6. Latency and memory claims are unsupported by the measurements}

\classification{Definite systems-reporting insufficiency}

\findinglabel{Location.} Experimental Setup; Tables~\texttt{tab:baselines} and \texttt{tab:equal\_alpha\_ablation}; Abstract; Conclusion.

\findinglabel{Problem.} Every non-RAG condition, including 12 ablation rows, negative steering, and the 20-vector placebo aggregate, is reported as exactly 7.32~s.  No run count, variance, warm-up, synchronization, batching, GPU count, timing boundary, prompt/output-length control, or peak-memory measurement is supplied.  The paper therefore cannot establish literal ``0 added generation latency overhead'' or a ``zero-memory advantage.''  Steering still stores a vector and hook state and uses model activations; it avoids a retrieval index and additional context tokens, which is a narrower claim.  It is also unclear whether the reported RAG time includes retrieval or only model execution.

\findinglabel{Why it matters.} Identical rounded means across many interventions can conceal measurement reuse or insufficient resolution.  The 87.1\% latency comparison is not interpretable without a common protocol and uncertainty.

\findinglabel{Suggested fix.} Run synchronized repeated timings after warm-up under fixed hardware, batch size, prompts, and decoding settings.  Report mean, dispersion, number of trials, prefill/decode/retrieval components, throughput, and peak CPU/GPU memory.  Replace ``zero overhead'' with ``no measurable overhead under this protocol'' only if an equivalence interval supports a prespecified negligible margin.  Replace ``zero-memory'' with ``no retrieval-index or additional context-token overhead.''

\subsection*{7. The ablation's small schedule rankings have no uncertainty and retain dose ambiguity}

\classification{Unsupported comparative claim and terminology error}

\findinglabel{Location.} contribution~5; cumulative-dose equations; Controlled Factorial Ablation Matrix; Discussion.

\findinglabel{Problem.} The manuscript correctly acknowledges that equal-$\alpha_0$ schedules have unequal $D_1$ and that the matched-$D_1$ row does not match temporal placement or $D_2$.  Nevertheless, the table emphasizes raw winners from differences as small as 0.0005--0.0018 in BERTScore without paired intervals or multiplicity handling.  The matched quantity is repeatedly called ``energy,'' although it is an $L_1$-type cumulative magnitude.  Moreover, the text defines a cumulative \emph{absolute} dose but writes
\[
D_{\text{continuous}}=T\alpha_0,
\qquad
D_{\text{decay}}=\frac{K+1}{2}\alpha_0,
\]
which become negative for the explicitly evaluated control $\alpha_0=-18$.  Absolute dose requires $|\alpha_0|$.

\findinglabel{Why it matters.} Readers may interpret numerical rank order as a reliable schedule advantage, and the sign mismatch makes the dose definition internally incorrect when applied to negative steering.

\findinglabel{Suggested fix.} Rename the control ``matched cumulative $L_1$ magnitude'' or ``matched $D_1$ dose,'' use absolute values in absolute-dose formulas, and reserve ``energy'' for a squared quantity such as $D_2$.  Report paired schedule contrasts with 95\% CIs and an exploratory multiplicity policy.  Avoid selecting a universal best schedule when ROUGE-L, BERTScore, and Rep-4 prefer different rows.

\subsection*{8. Vector extraction and layer-selection evidence do not establish a transferable factual direction}

\classification{Likely representation confound and unsupported mechanism inference}

\findinglabel{Location.} Contrastive Vector Estimation; Early-Stopping Steering Hook; Layer-Wise Subspace Probing.

\findinglabel{Problem.} The vector is estimated from the final token of the complete sequence $x_i\oplus y_i^{\pm}$, after the model has observed the entire correct or synthetic answer, but it is injected into the first 16 generation steps.  Final tokens may encode answer length, punctuation, templates, or synthetic-error artifacts rather than a context-invariant factual property.  The layer experiment is a steering-layer sweep, not a reported linear probe or subspace analysis.  Only the range 0.7040--0.7140 is given, with no per-layer results or uncertainty, yet the paper concludes that truthfulness representations are distributed across layers and that robustness is confirmed.

\findinglabel{Why it matters.} An intervention can improve the proxy without supporting the proposed representation-level interpretation.  The temporal mismatch also creates an implementation-facing generalization assumption.

\findinglabel{Suggested fix.} Compare vectors extracted from prompt-only states, answer-initial tokens, multiple token positions, and pooled response representations.  Test independently authored errors and source/template-disjoint data.  Report every layer's validation result and uncertainty, describe the experiment as a steering-layer sweep, and avoid conclusions about distributed truthfulness without direct probing or representation analysis.

\subsection*{9. Dataset provenance and leakage controls remain inadequate for a clinical benchmark}

\classification{Implementation ambiguity and likely leakage risk}

\findinglabel{Location.} ViHaluEval-Medical Benchmark Construction; Experimental Setup; contribution~1.

\findinglabel{Problem.} ``Expert-structured'' is undefined.  The paper does not state annotator qualifications, expert review fraction, agreement, adjudication, synthetic counter-answer generation procedure, quality-control rules, deduplication, source identifiers, or licensing.  A question-disjoint split does not ensure that questions based on the same drug/formulary entry or the same perturbation template are separated.  The selection procedure and seed for the 500-question subset are also not reported.

\findinglabel{Why it matters.} Source-entry or template leakage can inflate vector estimation, validation selection, and test reference preference.  Clinical claims require auditable provenance and safety checking.

\findinglabel{Suggested fix.} Add a dataset card and construction algorithm.  Split by drug/source entry and perturbation template before generating answer variants, publish split manifests and the evaluation-subset seed, report deduplication results, and document expert review and adjudication.  Obtain blinded clinical validation for a representative sample or all primary-test records.

\subsection*{10. Statistical, metric, and hook specifications remain incomplete}

\classification{Reproducibility limitation}

\findinglabel{Location.} BERTScore criterion; Experimental Setup; bounded-generation statistics; implementation description.

\findinglabel{Problem.} The 95\% interval $[+1.37,+7.03]$~pp is numerically consistent with a paired Wald-type interval from the displayed McNemar table, but its construction is not stated.  The Wilcoxon test is paired with a bootstrap interval for a mean difference without reporting the mean point estimate, Wilcoxon statistic, zero-difference handling, bootstrap procedure, or multiplicity policy.  Ties are said to be neither class, but their counts and their collapse into the McNemar ``non-preferred'' category are not explicit.  Rep-4 has no formula or tokenizer; Vietnamese ROUGE tokenization is unspecified; BERTScore omits package version, IDF, rescaling, and tie tolerance.  The hook omits the exact Qwen module, layer-index convention, modified tensor positions, prompt-forward behavior, KV-cache handling, dtype/device, and removal.  Software versions, seeds, GPU count, and complete generation arguments are also missing.

\findinglabel{Why it matters.} These choices affect outputs, effect estimates, and timing and prevent independent reproduction even when the displayed arithmetic is internally consistent.

\findinglabel{Suggested fix.} State all statistical estimators and software calls, report tie counts and paired effect estimates, and distinguish the Wilcoxon estimand from the bootstrap mean estimand.  Provide metric formulas/configurations, complete decoding arguments, environment versions, seeds, hardware count, and minimal executable hook pseudocode.  Release per-question outputs and analysis code.

\section*{Equation, notation, sign, branch, and dimensional audit}

The sign convention $\mu_l^+-\mu_l^-$ followed by addition of $+\alpha(t)v_{\text{steer}}^{(l)}$ is consistent with the intended positive direction, and unit normalization makes the vector dimensionally compatible with the hidden state when $\alpha(t)$ is interpreted as an activation-scale coefficient.  The linear schedule and the positive-$\alpha_0$ sums are arithmetically correct, including the final coefficient $\alpha(K)=\alpha_0/K$.  No inverse-trigonometric functions, coordinate transformations, or branch/quadrant choices occur.

The concrete audit findings are: the ``absolute dose'' formulas omit $|\alpha_0|$ even though negative steering is evaluated; $D_1$ is mislabeled as energy; $N$ denotes training examples, evaluation questions, category sizes, and random-vector count in different places; $v_{\mathrm{rand}}^{(1..20)}$ is not standard mathematical indexing; and ``truthfulness,'' ``factual preference,'' and ``reference preference'' drift between a representation interpretation and an observed proxy.  Define distinct symbols such as $N_{\mathrm{train}}$, $N_{\mathrm{test}}$, and $R$ random directions, and write $\{v_{\mathrm{rand}}^{(r)}\}_{r=1}^{R}$.

\section*{Meaningful editorial, compilation, and reference findings}

\subsection*{The source compiles, but two tables exceed the IEEE column width}

\findinglabel{Location.} Tables~\texttt{tab:long\_gen} and \texttt{tab:mcnemar}; clean pdfLaTeX log.

\findinglabel{Problem.} After two clean pdfLaTeX passes, the manuscript builds to a five-page PDF, but the log reports overfull boxes of approximately 29.75~pt and 5.54~pt for these tables.  Table~\texttt{tab:clinical\_safety} also declares six tabular columns with \verb|{lcrrrr}| while supplying five fields per row.

\findinglabel{Why it matters.} Content can extend into the neighboring IEEE column, and the extra alignment column makes the source inconsistent with the displayed table structure.

\findinglabel{Suggested fix.} Reduce or wrap the natural-completion and McNemar tables, or use a column-width resize only after checking readability.  Change the category table declaration to five columns (for example, \verb|{lcrrr}|), rebuild twice, and inspect the final PDF and log.

\subsection*{Several citations do not directly support the attached claims}

\findinglabel{Location.} Introduction; bibliography entries \texttt{ref10}, \texttt{ref11}, and \texttt{ref12}.

\findinglabel{Problem.} \texttt{ref10} is the LoRA paper rather than direct evidence for the manuscript's broad SFT hallucination-mitigation claim.  \texttt{ref11} studies robustness to irrelevant retrieved context but does not establish that models ignore evidence because ``internal activation pathways favor non-truthful generation.''  \texttt{ref12} establishes catastrophic forgetting generally, not the stated loss pattern for this medical SFT setting.

\findinglabel{Why it matters.} Mechanistic and comparative motivation should be separated from what the cited work actually demonstrates.

\findinglabel{Suggested fix.} Add direct evidence for hallucination-focused SFT, domain-adaptation forgetting in language models, and evidence-use failures in RAG.  Rewrite the internal-activation explanation as a hypothesis unless it is measured or directly supported.

\section*{Proofing sweep}

The bundled mechanical scan was run on both the current source and PDF.  Its only repeated hit was the notation \verb|v_{\text{rand}}^{(1..20)}|; the two periods are not prose punctuation, but the indexing should still be corrected.  A bounded manual check also found the following high-confidence defects:

\begin{itemize}
    \item \textbf{Directional-control notation:} replace \verb|v_{\text{rand}}^{(1..20)}| with $\{v_{\mathrm{rand}}^{(r)}\}_{r=1}^{20}$ or ``$v_{\mathrm{rand}}^{(1)},\ldots,v_{\mathrm{rand}}^{(20)}$.''
    \item \textbf{Placebo table:} remove ``283/500'' from the aggregate random-vector row; report the mean, SD, range, and number of directions instead.
    \item \textbf{Placebo and RAG prose:} replace ``a --10.20~pp drop'' and ``a --4.40~pp drop'' with ``a 10.20~pp drop'' and ``a 4.40~pp drop,'' or write signed changes as $-10.20$ and $-4.40$~pp.
    \item \textbf{Terminology:} replace ``un-oracle'' with ``non-oracle'' and ``Matched Energy Control'' with ``Matched $D_1$ Control'' or ``Matched Cumulative-Magnitude Control.''
    \item \textbf{Table notation:} define whether $56.50\%\pm5.11\%$ is mean $\pm$ sample SD across random directions; do not mix an aggregate mean with a single raw count.
\end{itemize}

The bibliography was checked for duplicated commas, malformed quotation punctuation, and obvious citation-format glitches; no additional mechanical defect was found.

\section*{Items requiring external verification}

\begin{itemize}
    \item Verify every formulary-derived reference answer and synthetic counter-answer against the official source, current clinical guidance, and qualified clinician review; also verify reuse permissions and the exact official title/metadata of the formulary.
    \item Validate the multilingual BERTScore configuration for Vietnamese clinical negation, quantities, units, contraindications, and interaction mechanisms before interpreting RefPref as factual alignment.
    \item Verify novelty against current work on token-dependent, scheduled, and dose-controlled activation steering; the present literature review does not establish that linear-decay early stopping is novel.
    \item Verify the characterization of Qwen2.5-7B-Instruct as a ``Small Language Model'' and the privacy/local-hospital deployment claims against explicit operational definitions, a threat model, and measured resource requirements.
    \item Verify the BM25 baseline against a standard Vietnamese retrieval configuration or stronger retrievers before making a general systems comparison with RAG.
\end{itemize}

\section*{Highest-priority revision sequence}

\begin{enumerate}
    \item Add clinician-adjudicated factual and harmful-error endpoints, or consistently narrow every claim to BERTScore reference preference.
    \item Test continuous, cutoff, decay, and dose-matched schedules under natural completion with paired uncertainty and direct repetition/activation diagnostics.
    \item Reanalyze and fully report the 20 random directions, negative steering, and BM25 outputs using paired or hierarchical uncertainty on identical questions.
    \item Complete the primary 800-token paired statistics and report Rep-4, lengths, finish reasons, and category effects.
    \item Document and strengthen the BM25 pipeline, including oracle context, Recall@$k$, retrieval-conditioned results, and component-wise latency/memory.
    \item Add dataset provenance, source/template-disjoint splits, clinical review, vector-extraction controls, and complete metric/hook implementation details.
    \item Correct the $D_1$ terminology/sign formulas, aggregate-table errors, overfull tables, indexing notation, double-minus copy, and citation fit.
\end{enumerate}

\end{document}